\addcontentsline{toc}{section}{СПИСОК ИСПОЛЬЗОВАННЫХ ИСТОЧНИКОВ}

\begin{thebibliography}{9}
	%Ссылка есть 
	\bibitem{zakon} Российская Федерация. Законы. О применении контрольно-кассовой техники при осуществлении расчетов в Российской Федерации: Федеральный закон № 54-ФЗ от 22 мая 2003 г. - Текст :
	электронный // КонсультантПлюс: справочно-правовая система. – URL:https://www.consultant.ru/document/cons\_doc\_LAW\_42359/ (дата обращения: 02.10.2024).
	\bibitem{zzakon} Федеральная налоговая служба Российской Федерации. Приказ ФНС России от 31 декабря 2018 г. № ММВ-7-20/750@ «Об утверждении порядка применения контрольно-кассовой техники» : приказ ФНС РФ от 31.12.2018 № ММВ-7-20/750@ – Текст : электронный // Федеральная налоговая служба официальный сайт – URL:https://www.nalog.gov.ru/rn77/about\_fts/docs/11715492/(дата обращения: 02.10.2024).
    \bibitem{tkitner} Букунов С. В. Разработка приложений с графическим пользовательским интерфейсом на языке Python : учебное пособие для СПО / С. В. Букунов, О. В. Букунова. — Санкт-Петербург : Лань, 2023. — 88 с. - ISBN 978-5-507-45192-0 - Текст: непосредственный.
    \bibitem{postgrestermin} Джуба	С.,	Волков	А. Изучаем PostgreSQL 10 / пер. с анг. А. А. Слинкина. – М.: ДМК Пресс, 2019. – 400 с. - ISBN 978-5-97060-643-8 - Текст: непосредственный.
    \bibitem{bisness}Кулагин В., Сухаревски А., Мефферт Ю. Digital@Scale : Настольная книга по цифровизации бизнеса / Владимир Кулагин, Александр Сухаревски, Юрген Мефферт. М. : Интеллектуальная Литература, 2019 — 293 с. - ISBN 978-5-6042320-7-1 -  Текст: непосредственный.
    \bibitem{guipython}	Машнин Т. Создание настольных Python приложений с графическим интерфейсом пользователя / Машин Т. – Москва: ЛитРес: Самиздат, 2021. – 131 с. – ISBN 978-5-532-96908-7. – Текст~: непосредственный.
    \bibitem{uml} Моисеев А.Н., Литовченко М.И. Основы языка UML : учеб. пособие. – Томск : Издательство Томского государственного университета, 2023. – 96 с. - ISBN 978-5-907722-06-4 - Текст: непосредственный.
    \bibitem{components} Новиков, Ф. А. Учебно-методическое пособие по дисциплине «Анализ и проектирование на UML» : учебно-методическое пособие / Ф. А. Новиков. — Санкт-Петербург : НИУ ИТМО, 2007. — 286 с. — Текст : электронный.
    \bibitem{python} Россум Г. , Ф.Л.Дж. Дрейк, Д.С. Откидач, [и др.]. Язык программирования Python. / 2001 — 454 c. - ISBN 978-9-667-39354-0  — Текст : непосредственный.
    \bibitem{statiay} Скоробогатов М.В., Минченко Л.В. Внедрение инструментов цифровизации в сфере общественного питания / Научный журнал НИУ ИТМО. Серия "Экономика и экологический менежджмент" - 2023. - №1 - с.108-116.
    \bibitem{postgres} Ткачев О. А.  PostgreSQL: SQL + PL/pgSQL для тех, кто хочет стать профессионалом — СПб.: Издательство Наука и Техника, 2024. — 480 с.- ISBN 978-5-907592-32-2 - Текст: непосредственный.
    \bibitem{siteallsys} Системы управления рестораном: 2025: сайт / Подбор и внедрение лучших IT-решений для работы и бизнеса. - URL: https://picktech.ru/catalog/restaurant-management-software/(дата обращения 02.05.2025). - Текст: электронный.
	\bibitem{ikko} IikoOficce - Руководство пользователя / iikoRMS (версия 7.3). Руководство пользователя iikoOffice. - 2020. - 810с - Текст: Электронный.
	\bibitem{keeper} Профессиональная система R-KEEPER для ресторанов Руководство кассириа, бармена, официанта - Москва - 2017. - 247 с. - Текст: электронный. 
	\bibitem{onec} Использование конфигурации «Удобное решение: Ресторан 1.0» - Москва - 2014. - 49 с. - Текст: электронный.
	\bibitem{siteanaliz} Автоматизация ресторана и кафе: системы и программы для работы, управления, учёта: сайт / Компания Клеверенс - URL:https://www.cleverence.ru/articles/auto-busines/avtomatizatsiya-restorana-i-kafe-sistemy-i-programmy-dlya-raboty-upravleniya-uchyeta/ (дата обращения: 20.09.2024). – Текст: электронный.
	\bibitem{gosty} Кисилевич Т. И. Автоматизация деятельности предприятий общественного питания: теория и опыт / Т. И. Кисилевич, Я. Ю. Митрюшкин, Г. Р. Хвистани // Инновационное развитие экономики. – 2018. – № 6-1 (48). – С. 161–166.
	\bibitem{statoptim} Шихатов П. И. Автоматизация учетных процессов на предприятиях общественного питания и ее роль в формировании современной модели управленческого учета [Текст] / П. И. Шихатов // Вестник МИЭП. – 2017. – № 2.– С. 21–24. 
	\bibitem{sitkin} Ситкин, В. А. Автоматизация ресторана и кассовая техника: учебное пособие / В. А. Ситкин, Н. В. Химич. – Москва : Сервис ККМ, 2010.– 132 с.  - Текст: электронный. 
	\bibitem{gost}ГОСТ Р 50647-94. Общественное питание. Термины и определения : национальный стандарт Российской Федерации : дата введения 1995-01-01 / разработан Федеральным агентством по техническому регулированию и метрологии. – Москва : Издательство стандартов, 1995. – 10 с. – Текст : непосредственный.
\end{thebibliography}

